\documentclass[11pt]{article}

\usepackage[a4paper,margin=1in]{geometry}
\usepackage{amsmath}
\usepackage{graphicx}
\usepackage{booktabs}
\usepackage{array}
\usepackage{listings}
\usepackage{xcolor}
\usepackage{hyperref}
\usepackage{caption}
\usepackage{float}
\usepackage{enumitem}
\usepackage{fancyhdr}

\hypersetup{colorlinks=true, linkcolor=blue, citecolor=blue, urlcolor=blue}

\definecolor{codebg}{rgb}{0.96,0.96,0.96}
\lstset{
  backgroundcolor=\color{codebg},
  basicstyle=\ttfamily\small,
  breaklines=true,
  frame=single,
  numbers=left,
  numberstyle=\tiny,
  keywordstyle=\color{blue!70!black}\bfseries,
  commentstyle=\color{green!40!black},
  showstringspaces=false,
  tabsize=4
}

\pagestyle{fancy}
\fancyhf{}
\rhead{Lab 3: Process Scheduling}
\lhead{Skanda P, Narapogu Suhas}
\cfoot{\thepage}

\title{\textbf{Laboratory 3: Process Scheduling}\\[4pt]
\large Simulation and Comparative Analysis of CPU Scheduling Algorithms}
\author{Skanda P \; (CS24BT045) \and Narapogu Suhas \; (CS24BT019)}
\date{}

\begin{document}

\maketitle
\thispagestyle{fancy}

% ============================================================
\section{Introduction}
% ============================================================
The goal of this assignment is to build a simulator for CPU scheduling
algorithms and use it to study how different scheduling policies affect
process turnaround time. As specified in the assignment, the simulator reads
a workload description file in which each line represents one process:

\begin{center}
\texttt{<arrival-time> <cpu-burst-1> <io-burst-1> <cpu-burst-2> <io-burst-2> ... -1}
\end{center}

and is invoked from the command line as

\begin{center}
\texttt{./simulator <algorithm> <path-to-workload-file>}
\end{center}

(with an optional numeric parameter -- time quantum for Round Robin, or
boost interval for MLFQ -- inserted between the algorithm name and the
workload file path).

The simulator produces two things: (1) the resulting \emph{schedule}, in the
format given in \texttt{schedule\_format.txt}, and (2) summary statistics:
the average turnaround time, the maximum turnaround time, and the
simulator's run time (the total CPU-burst time, excluding I/O).

Part I of the assignment implements FIFO, Round Robin, and MLFQ for a
single processor. Part II re-evaluates the same three algorithms with two
processors sharing a common set of ready queues.

% ============================================================
\section{Design of the Simulator}
% ============================================================

\subsection{Process representation}
Every process is represented by a single \texttt{struct} that stores both the
static description read from the workload file and the dynamic state needed
during simulation:

\begin{lstlisting}[language=C++, caption={Process data structure (\texttt{src/simulator\_part1.cpp})}]
struct Process {
    int pid;
    int arrival_time;
    vector<int> bursts;   // even indices: CPU bursts, odd indices: I/O bursts
    int curr_b_idx;       // index of the burst currently being served
    int remaining_b_time; // time left in the current burst
    int completion_time;  // time at which the process finishes entirely
    int queue_level;      // MLFQ priority level: 0 (highest) .. 2 (lowest)
};
\end{lstlisting}

Storing the CPU and I/O bursts of a process in a single alternating vector
(\texttt{bursts[0]} = CPU, \texttt{bursts[1]} = I/O, \texttt{bursts[2]} =
CPU, \ldots) keeps the description compact and lets the simulator move
between phases simply by incrementing \texttt{curr\_b\_idx}.

\subsection{Queues}
The simulator maintains:
\begin{itemize}[nosep]
  \item a single FIFO \texttt{ready\_queue} used by both the FIFO and Round
        Robin algorithms,
  \item three FIFO queues \texttt{mlfq\_queues[0..2]} for MLFQ, with
        \texttt{mlfq\_queues[0]} the highest-priority queue, and
  \item a list \texttt{io\_waiters} of processes currently performing I/O.
\end{itemize}

\subsection{Discrete-time event loop}
The simulator advances time one unit ("tick") at a time. In every tick it:
\begin{enumerate}[nosep]
  \item admits any processes whose arrival time equals the current time
        into the ready queue (or \texttt{mlfq\_queues[0]} for MLFQ);
  \item applies a periodic priority boost for MLFQ if the boost interval has
        elapsed, moving every process in \texttt{mlfq\_queues[1..2]}, every
        process doing I/O, and the currently running process back to queue
        level 0;
  \item dispatches an idle CPU from the highest non-empty queue if one is
        available;
  \item executes the running process(es) and all processes performing I/O
        for one time unit;
  \item checks for I/O completion (moving processes back into the ready
        queues) and for CPU burst completion or time-quantum expiry
        (recording a schedule entry and either retiring the process,
        moving it to I/O, or re-queuing it, possibly at a lower MLFQ
        level).
\end{enumerate}
This tick-based design keeps FIFO, Round Robin, and MLFQ under a single,
uniform simulation loop: FIFO is simply Round Robin with an infinite
quantum (quantum $=0$ disables preemption), and MLFQ reuses the same
dispatch/preemption logic with three queues and a quantum of 2 for every
level.

\subsection{Part II: two processors}
For Part II (\texttt{src/simulator\_part2.cpp}), the ready queue(s) and the
I/O-waiting pool remain \emph{shared} across processors, but each of the two
processors keeps its own currently-running process, its own quantum
counter, and its own schedule log. On every tick, both idle processors are
allowed to dequeue a process (dispatch phase) before either processor
executes, and both processors then advance their assigned process by one
time unit in the same tick, faithfully modeling true parallel execution.
The output schedule is printed once per processor (\texttt{CPU0},
\texttt{CPU1}), matching the format required by
\texttt{schedule\_format.txt}.

\subsection{Command line interface}
Both simulators are invoked as
\begin{lstlisting}
./sim_part1 FIFO <workload-file>
./sim_part1 RR <quantum> <workload-file>
./sim_part1 MLFQ <boost-interval> <workload-file>
\end{lstlisting}
(and analogously \texttt{./sim\_part2} for two processors). A boost interval
of \texttt{0} for MLFQ disables the periodic priority boost.

% ============================================================
\section{Part I: Single-Processor Scheduling}
% ============================================================

All experiments below use the five workload description files supplied with
the assignment, located in \texttt{src/test\_cases/} and named
\texttt{process1.txt} through \texttt{process5.txt}. We refer to them as
Workload 1 -- Workload 5. Workload
4 (three processes, several CPU/I/O cycles) and Workload 5 (six processes,
four of which arrive together with identical burst patterns) are used for
the quantum and boost studies below, since they are the most representative
of realistic multiprogrammed workloads; Workloads 1--3 are simpler
sanity-check inputs (pure-CPU jobs, and short/cheap I/O jobs respectively).

\subsection{FIFO}
Table~\ref{tab:fifo} shows the turnaround-time statistics obtained by running
FIFO on all five workloads.

\begin{table}[H]
\centering
\caption{FIFO results, single processor}
\label{tab:fifo}
\begin{tabular}{lccc}
\toprule
Workload & Avg.\ Turnaround & Max.\ Turnaround & Run Time \\
\midrule
1 & 363.33 & 590 & 600 \\
2 & 525.00 & 600 & 600 \\
3 &  14.50 &  16 &  12 \\
4 & 123.00 & 163 & 150 \\
5 & 189.00 & 225 & 240 \\
\bottomrule
\end{tabular}
\end{table}

Because FIFO is non-preemptive, once a process starts running it keeps the
CPU until it either finishes a CPU burst (and moves to I/O) or terminates.
This is visible in the sample schedule for Workload 4 in
Table~\ref{tab:sample-schedule}: \texttt{P1} runs uninterrupted for its
entire first burst (time 1--30) before \texttt{P2} is even allowed to start,
even though \texttt{P2} and \texttt{P3} arrived earlier or at the same time
as some of \texttt{P1}'s later bursts. Consequently, FIFO strongly favors
processes that are scheduled early and can starve short jobs that arrive
slightly later -- \texttt{P3} in Workload 4, despite having very short CPU
bursts, ends up with the largest turnaround time (163) simply because it is
queued behind the two longer processes.

\begin{table}[H]
\centering
\caption{Sample schedule excerpt: FIFO on Workload 4 (\texttt{P<pid>,<burst\#>\ \ start\ \ end})}
\label{tab:sample-schedule}
\begin{tabular}{lccc}
\toprule
Entry & Start & End & \\
\midrule
P1,1 & 1   & 30  \\
P2,1 & 31  & 40  \\
P3,1 & 41  & 43  \\
P1,2 & 44  & 73  \\
P2,2 & 74  & 83  \\
P3,2 & 84  & 86  \\
\ldots & & & (P3 is repeatedly delayed behind P1/P2) \\
P3,10& 176 & 178 \\
\bottomrule
\end{tabular}
\end{table}

\subsection{Round Robin}
We swept the time quantum over $\{1, 2, 4, 8\}$ time units and measured the
average turnaround time for Workload 4 and Workload 5
(Figure~\ref{fig:rr1cpu}, Table~\ref{tab:rr1cpu}).

\begin{table}[H]
\centering
\caption{Round Robin: average turnaround time vs.\ quantum, single processor}
\label{tab:rr1cpu}
\begin{tabular}{lcccc}
\toprule
Quantum & 1 & 2 & 4 & 8 \\
\midrule
Workload 4 & 130.00 & 131.33 & 131.67 & 128.33 \\
Workload 5 & 203.83 & 215.67 & 207.17 & 205.67 \\
\bottomrule
\end{tabular}
\end{table}

\begin{figure}[H]
  \centering
  \includegraphics[width=0.7\textwidth]{figs/rr_quantum_1cpu.png}
  \caption{Average turnaround time as a function of the Round Robin time
  quantum, single processor.}
  \label{fig:rr1cpu}
\end{figure}

\paragraph{Observations.} The turnaround time does not vary monotonically
with the quantum; instead it is fairly flat and workload-dependent. A very
small quantum ($q=1$) gives the best average turnaround for both workloads
here because it lets short/interactive bursts interleave tightly with
longer ones, closely approximating processor sharing, but it does so at the
cost of a much larger number of context switches (visible as many more,
shorter entries in the schedule than under FIFO). A very large quantum
($q=8$) starts to behave like FIFO for bursts shorter than the quantum,
which is why its performance on Workload 4 is close to the $q=1$ case for
this particular process mix. Intermediate quanta ($q=2,4$) are
occasionally slightly worse here because they interact less favorably with
the specific burst lengths in the workload (e.g.\ a 30-unit burst split
into chunks of 4 versus chunks of 8 changes exactly which process is
waiting when I/O bursts complete). This illustrates the classical
trade-off: very small quanta approach fairness/interactivity at the cost of
scheduling overhead, while very large quanta approach FIFO behaviour and
can reintroduce the same starvation-of-short-jobs effect that plain FIFO
exhibits.

\subsection{Multi-Level Feedback Queue (MLFQ)}
MLFQ was implemented with three queues $Q_0$ (highest priority) to $Q_2$
(lowest), each with a time quantum of 2. A process enters $Q_0$ on arrival;
if it uses its full quantum without finishing its burst it is demoted one
level (to a minimum of $Q_2$), and it is returned to whichever level it was
at after finishing an I/O burst. We compared no periodic boost against a
boost every 20 time units, where the boost resets every process (whether
in a queue, running, or doing I/O) back to $Q_0$.

\begin{table}[H]
\centering
\caption{MLFQ: effect of periodic priority boost, single processor}
\label{tab:mlfq1cpu}
\begin{tabular}{lcc}
\toprule
Workload & No boost & Boost every 20 \\
\midrule
4 & 135.33 & 131.67 \\
5 & 212.83 & 206.17 \\
\bottomrule
\end{tabular}
\end{table}

\begin{figure}[H]
  \centering
  \includegraphics[width=0.65\textwidth]{figs/mlfq_boost_1cpu.png}
  \caption{Average turnaround time with and without periodic priority
  boosting, single processor.}
  \label{fig:mlfq1cpu}
\end{figure}

\paragraph{Observations.} Without a boost, a long-running, CPU-bound
process is quickly demoted to $Q_2$ and, if new processes keep arriving and
occupying $Q_0$/$Q_1$, it can be starved for a long time -- exactly the
classical MLFQ starvation problem. Periodically boosting every process back
to $Q_0$ every 20 time units gives previously-demoted processes another
chance to run at high priority, which reduces both the average and, in our
experiments, the tail turnaround times slightly. The effect is modest for
these particular (short) workloads because the total run length is small
relative to the boost interval, but the direction of the effect -- boosting
reduces average turnaround time by reducing starvation of long-running
processes -- matches the expected behaviour of MLFQ described by Arpaci-Dusseau
and Arpaci-Dusseau's ``rule of thumb'': periodic boosting guarantees
CPU-bound processes are not permanently starved.

% ============================================================
\section{Part II: Two-Processor Scheduling}
% ============================================================

For Part II, the same three algorithms were re-evaluated with the queues
shared between two processors (\texttt{sim\_part2}). Both processors dequeue
independently from the shared ready queue(s) but execute in the same
simulated tick, so a workload with enough parallel-CPU-burst opportunities
should see close to a $2\times$ reduction in turnaround time (limited by
inherently sequential dependencies, e.g.\ a single process cannot use both
processors at once, and I/O bursts remain serial per process).

\subsection{FIFO}
\begin{table}[H]
\centering
\caption{FIFO results, two processors}
\label{tab:fifo2}
\begin{tabular}{lccc}
\toprule
Workload & Avg.\ Turnaround & Max.\ Turnaround & Run Time \\
\midrule
1 & 230.00 & 390 & 600 \\
2 & 275.00 & 300 & 600 \\
3 &  13.50 &  14 &  12 \\
4 &  81.33 &  93 & 150 \\
5 & 104.33 & 119 & 240 \\
\bottomrule
\end{tabular}
\end{table}

\subsection{Round Robin}
\begin{table}[H]
\centering
\caption{Round Robin: average turnaround time vs.\ quantum, two processors}
\label{tab:rr2cpu}
\begin{tabular}{lcccc}
\toprule
Quantum & 1 & 2 & 4 & 8 \\
\midrule
Workload 4 & 81.67 & 79.67 & 81.67 & 81.67 \\
Workload 5 & 107.50 & 107.33 & 106.83 & 105.67 \\
\bottomrule
\end{tabular}
\end{table}

\begin{figure}[H]
  \centering
  \includegraphics[width=0.7\textwidth]{figs/rr_quantum_2cpu.png}
  \caption{Average turnaround time as a function of the Round Robin time
  quantum, two processors.}
  \label{fig:rr2cpu}
\end{figure}

\paragraph{Observations.} With two processors the sensitivity of average
turnaround time to the choice of quantum shrinks considerably (the curves
in Figure~\ref{fig:rr2cpu} are much flatter than in
Figure~\ref{fig:rr1cpu}). With extra CPU capacity, processes spend far less
time waiting in the ready queue in the first place, so the quantum -- which
mainly governs how ready-queue waiting time is distributed among
runnable processes -- has less to redistribute.

\subsection{MLFQ}
\begin{table}[H]
\centering
\caption{MLFQ: effect of periodic priority boost, two processors}
\label{tab:mlfq2cpu}
\begin{tabular}{lcc}
\toprule
Workload & No boost & Boost every 20 \\
\midrule
4 & 80.67 & 80.33 \\
5 & 107.50 & 110.17 \\
\bottomrule
\end{tabular}
\end{table}

\paragraph{Observations.} With two processors, both queues drain quickly
enough that processes rarely accumulate the amount of ready-queue waiting
time needed to be demoted far enough for starvation to become a real risk;
consequently periodic boosting has little to no benefit, and for Workload 5
it is even marginally worse, because the boost briefly promotes
long-running processes ahead of newly-demoted short ones, adding a small
amount of extra waiting for the latter. This is consistent with the
general expectation that priority-boost tuning matters most when the
system is CPU-constrained.

\subsection{Comparing one processor against two processors}

\begin{figure}[H]
  \centering
  \includegraphics[width=0.75\textwidth]{figs/compare_1v2_wl4.png}
  \caption{Workload 4: average turnaround time, 1 CPU vs.\ 2 CPUs, for each
  algorithm (using the best-performing quantum/boost setting from Part I).}
  \label{fig:compare4}
\end{figure}

\begin{figure}[H]
  \centering
  \includegraphics[width=0.75\textwidth]{figs/compare_1v2_wl5.png}
  \caption{Workload 5: average turnaround time, 1 CPU vs.\ 2 CPUs, for each
  algorithm.}
  \label{fig:compare5}
\end{figure}

Across every algorithm and workload we tested, adding a second processor
roughly halves the average turnaround time (e.g.\ FIFO on Workload 4 drops
from 123.0 to 81.3, and on Workload 5 from 189.0 to 104.3). The reduction
is largest for FIFO, since FIFO's main weakness -- a process blocking the
CPU while others wait -- is directly alleviated by a second CPU that can
serve a waiting process concurrently. Round Robin and MLFQ, which already
mitigate that weakness through preemption on a single CPU, show a smaller
\emph{relative} improvement, though the improvement in \emph{absolute}
turnaround time is still substantial. Interestingly, the ranking of
algorithms by average turnaround time is not always preserved when going
from one to two processors: MLFQ and Round Robin are the best performers on
one CPU, but with two CPUs FIFO becomes competitive with (and on Workload 5,
slightly better than) MLFQ, since with ample CPU capacity the preemption
overhead of RR/MLFQ (extra context switches without a corresponding
reduction in queue waiting) yields no further benefit.

% ============================================================
\section{Summary of Findings}
% ============================================================
\begin{itemize}
  \item \textbf{FIFO} is simple and has zero preemption overhead, but is
        highly sensitive to arrival order: a long process can force much
        shorter processes that arrive shortly after it to wait
        disproportionately long, inflating both average and maximum
        turnaround time.
  \item \textbf{Round Robin} evens out waiting time between processes by
        time-slicing the CPU. Very small quanta approach ideal processor
        sharing (best average turnaround in our experiments) at the cost of
        more context switches; very large quanta degenerate toward FIFO.
        The best quantum is workload-dependent, so it is worth tuning quantum
        to the typical CPU-burst length in the system.
  \item \textbf{MLFQ} adapts to each process's observed behaviour (I/O-bound
        processes stay near the top, CPU-bound processes sink to the
        bottom) without any prior knowledge of burst lengths. Its main
        weakness is potential starvation of demoted, CPU-bound processes;
        a periodic priority boost (every 20 time units in our study) is a
        simple, effective way to fix this and gave a small but consistent
        reduction in average turnaround time on a single CPU.
  \item \textbf{Adding a second processor} improves every algorithm
        substantially by directly cutting ready-queue waiting time,
        and it shrinks the performance gap between the different
        scheduling policies, since with sufficient CPU capacity the
        scheduling policy matters less than raw availability of a free
        CPU.
\end{itemize}

% ============================================================
\section{Reproducing the Experiments}
% ============================================================
The complete source code, a \texttt{Makefile}, and the sample workload
files used in this report are included with the submission. To build and
run:
\begin{lstlisting}
# builds ./sim_part1 (1 CPU) and ./sim_part2 (2 CPUs)
make

./sim_part1 FIFO  src/test_cases/process4.txt
./sim_part1 RR  2 src/test_cases/process4.txt
./sim_part1 MLFQ 20 src/test_cases/process4.txt

./sim_part2 FIFO  src/test_cases/process4.txt
./sim_part2 RR  2 src/test_cases/process4.txt
./sim_part2 MLFQ 20 src/test_cases/process4.txt
\end{lstlisting}
Each invocation prints the schedule (per CPU) followed by the
\texttt{Final Results} block containing per-process turnaround time,
average turnaround time, maximum turnaround time, and simulator run time.

\end{document}
